% TOP_PROBLEM: {"id": 30006714, "title": "Abundance Conjecture for the minimal model program", "category_id": 5, "category": {"id": 5, "name": "algebraic_geometry", "display_name": "Algebraic Geometry", "description": "Geometric objects defined by polynomial equations.", "slug": "algebraic-geometry", "order_index": 5, "created_at": "2026-07-31T15:26:25.671Z"}, "status": "open"}
% FIRST_POSTED: 2026-09-25
% ENTRY_KIND: statement_only
% !TeX program = lualatex
% Definition and sources only; this file is not an LLM solution attempt.
\documentclass[11pt]{article}
\usepackage[margin=25mm]{geometry}
\usepackage{fontspec}
\setmainfont{Latin Modern Roman}
\usepackage{amsmath,amssymb,mathtools}
\usepackage{unicode-math}
\setmathfont{Latin Modern Math}
\usepackage{xurl,hyperref}
\hypersetup{colorlinks=true,urlcolor=blue}
\setcounter{secnumdepth}{0}
\setlength{\parindent}{0pt}
\setlength{\parskip}{5pt}
\setlength{\emergencystretch}{3em}
\begin{document}
\section*{\texorpdfstring{Abundance Conjecture for the minimal model program}{Abundance Conjecture for the minimal model program}}\label{30006714}
\subsection{Definitions and mathematical statement}
For a projective klt pair in the source's minimal-model setting, put $D=K_X+\Delta$. A divisor is nef if $D\cdot C\ge0$ for every curve $C$. In the usual rational-divisor formulation it is semiample if for some integer $m>0$, $mD$ is Cartier and its line bundle is generated by global sections. The conjecture is
\[
 D\text{ nef}\quad\Longrightarrow\quad \exists m>0:\mathcal O_X(mD)\text{ globally generated}.
\]
The rational-boundary convention is the standard one for this displayed statement; a real-boundary version requires the source's corresponding notion of semiampleness. No bigness hypothesis is added.
\par\smallskip\textit{Formulation and concept references:} \href{https://arxiv.org/abs/2406.18233}{[S1]}.

\subsection{Short English statement}
Must a nef canonical-plus-boundary divisor on a minimal klt variety eventually have enough sections to define a morphism?
\subsection{Sources}
\begin{itemize}
\item[S1]Vladimir Lazic, Metrics with minimal singularities and the Abundance conjecture (2024).
\url{https://arxiv.org/abs/2406.18233}.
\textit{Formulation source.}
\item[S2]Inequalities of Miyaoka-Yau type and Uniformisation of varieties of intermediate Kodaira Dimension.
\url{https://arxiv.org/abs/2601.15138}.
\textit{Further reference; not a proof endorsement.}
\item[S3]An approach to the abundance conjecture for Kaehler varieties via algebraic reduction.
\url{https://arxiv.org/abs/2604.03879}.
\textit{Further reference; not a proof endorsement.}
\item[S4]Two non-vanishing results concerning the anti-canonical bundle.
\url{https://www.cambridge.org/core/journals/nagoya-mathematical-journal/article/two-nonvanishing-results-concerning-the-anticanonical-bundle/1E68F2B88B818A7CBD47B96B7D700548}.
\textit{Further reference; not a proof endorsement.}
\end{itemize}
\end{document}
